\begin{resumo}
Este relatório documenta a segunda avaliação parcial do projeto Higrolar (equipe Estiagem, eq04), uma plataforma residencial de medição de temperatura e umidade relativa com ESP32, DHT11 e OLED SSD1306. O protótipo foi integrado a uma rede Wi-Fi de 2,4 GHz e ao broker EMQX Cloud por MQTT, com telemetria QoS 1, disponibilidade anunciada por estado retido e Last Will, configuração remota do período e armazenamento por assinante Python em CSV. Uma campanha física, realizada de 23/10/2026 às 14h27min27s até 24/10/2026 às 15h56min54s (UTC-3), produziu 3.096 mensagens, incluindo 3.082 telemetrias; a chave lógica (id, boot\_id, seq) revelou 3.080 medições únicas, duas duplicações e nenhuma lacuna interna nas sequências observadas. Foram registradas três inicializações, um horário de medição nulo antes da sincronização NTP e nenhuma falha de leitura do sensor registrada. A temperatura variou de 26,9 a 36,9~\si{\celsius}, e a umidade relativa, de 34\% a 73\%. Experimentos documentados envolveram queda deliberada de Wi-Fi, reinício do nó, interrupção do gravador, umidificação, refrigeração e operação simultânea de aparelhos. No período de 24/10 das 12h59 às 14h10, com ar-condicionado e umidificador ligados, a média observada foi de 30,14~\si{\celsius} e 59,28\% de umidade, sem inferência causal, pois o ambiente não foi controlado. A solução confirmou o caminho completo de dados até o disco e oferece registros temporais úteis à rotulagem no ciclo 3. Limitações são a resolução do DHT11, a fila volátil, a validação TLS não habilitada no firmware embarcado sem CA e a falta de medição instrumental do consumo do rádio.
\vspace{\onelineskip}
\noindent\textbf{Palavras-chave:} Internet das Coisas; ESP32; MQTT; EMQX Cloud; telemetria; qualidade de dados.
\end{resumo}
\begin{resumo}[Abstract]
\begin{otherlanguage*}{english}
This report presents the second stage of Higrolar, a residential temperature and relative-humidity monitoring system based on ESP32, DHT11 and an SSD1306 display. The device uses Wi-Fi and MQTT over an EMQX Cloud broker, publishing QoS 1 telemetry, retained presence status, Last Will messages and remotely configurable sampling periods. A physical campaign from October 23, 2026 at 14:27:27 to October 24, 2026 at 15:56:54 (UTC-3) recorded 3,096 MQTT messages, including 3,082 telemetry messages and 3,080 logically unique readings identified by (id, boot\_id, seq). Two logical duplicates, no detectable internal sequence gaps, three observed boot identifiers, and one missing acquisition timestamp were found. Recorded temperatures ranged from 26.9 to 36.9 degrees Celsius and relative humidity from 34 to 73 percent. Documented tests covered Wi-Fi disruption, recorder reconnection, reboot, humidification, air conditioning and simultaneous operation. From 12:59 to 14:10 on October 24, both devices were running, while the recorded mean was 30.14 degrees Celsius and 59.28 percent relative humidity; this is an observational association rather than causal evidence. The system delivered sensor data to a persistent CSV dataset. Limitations include sensor resolution, volatile local buffering, missing ESP32 server-certificate verification without a configured CA and lack of instrumented radio-current measurements.
\vspace{\onelineskip}
\noindent\textbf{Keywords:} Internet of Things; ESP32; MQTT; EMQX Cloud; telemetry; data quality.
\end{otherlanguage*}
\end{resumo}
